% Generated by GrindEQ Word-to-LaTeX 
\documentclass{article} %%% use \documentstyle for old LaTeX compilers

%\usepackage[english]{babel} %%% 'french', 'german', 'spanish', 'danish', etc.
\usepackage{amssymb}
\usepackage{amsmath}
%\usepackage{txfonts}
%\usepackage{mathdots}
%\usepackage[classicReIm]{kpfonts}
\usepackage{fancyhdr}
\usepackage{geometry}
\usepackage{lastpage}
\geometry{letterpaper,top=50pt,hmargin={20mm,20mm},headheight=15pt}
\usepackage[pdftex]{graphicx} %%% use 'pdftex' instead of 'dvips' for PDF output
\usepackage[utf8]{inputenc}

% You can include more LaTeX packages here 

\pagestyle{fancy}
\fancyhf{}
\chead{Day 21 reading -- typos corrected}
\lhead{E340 Nov. 24, 2021}
\rhead{Page \thepage/\pageref{LastPage}}

\begin{document}

Nov 24: fixed feedback variable:  they were denoted by $c_{pl}$, changed to $f_{pl}$

\section{Bathtub with feedbacks}


\textbf{Also see  Dessler 6.3 which covers the same material from a slightly different angle}


Remember our bathtub slide:

\begin{figure}[htbp]
    \centering
\includegraphics*[width=0.8 \textwidth]{bathtub.png} 
\caption{The bathtub}
  \label{fig:bathtub}
\end{figure}



 The stock and and flow equation for this system looks like: 
 \begin{equation}
   \label{ZEqnNum118824} 
\frac{dV}{dt} =I_{in} -I_{out} =I_{in} -kV 
\end{equation} 
Where, for the worksheet, V is in liters and t is in minutes, so k has to have units of per minute, or min${}^{-1}$.  It is more common to work with this version of \eqref{ZEqnNum118824}:
\begin{equation}
  \label{ZEqnNum264554} 
\frac{dV}{dt} =I_{in} -I_{out} =I_{in} -\frac{V}{\lambda }  
\end{equation} 
Where $\lambda=1/k$ is called the ``sensitivity'' and for this problem has units of minutes.  The bigger outflow drainpipe, the larger the value of k and the smaller the value of $\lambda$. Note that we can write down the final equilibrium volume for the tub by setting dV/dt=0 in \eqref{ZEqnNum264554}:


\begin{equation}
  V_{final} =\lambda I_{in} 
\end{equation}

In words, the smaller the outflow pipe, and the bigger the inflow rate, the higher the water level will have to be to reach equilibrium.

 To make this analogy work a little better for the climate case, let's assume that we start out in a state of equilibrium so that  we have:
 \begin{equation}
   \label{4)} 
V_{0} =\lambda I_{0}  
\end{equation} 
Then we increase the inflow to $I_{in}$ and the volume starts to increase.  If we write
$\Delta V= V - V_{0}$ and $\Delta I=I_{in} - I_{0}$, then we can add $V_{0}/\lambda$ and
subtract $I_{0}$ from the right hand side of \eqref{ZEqnNum264554} like this:

\begin{equation}
  \label{eq:addI0}
  \frac{dV}{dt}  =I_{in} -\frac{V}{\lambda }  - I_0 + \frac{V_0}{\lambda} = (I_{in} - I_{0}) - \frac{V - V_0}{\lambda }
\end{equation}
Which we can write as:
\begin{equation}
  \label{ZEqnNum448161} 
\frac{\Delta V}{dt} =\Delta I_{in} -\frac{\Delta V}{\lambda }  
\end{equation} 
(note that we've also added dV${}_{0}$/dt to the left-hand side of \eqref{ZEqnNum264554} which is the same as adding 0 since $V_{0}$ is a constant). 

  

 You should convince yourself by taking the derivative with respect to time and that the following equation is the solution to \eqref{ZEqnNum448161}:

 \begin{equation}
   \label{ZEqnNum190537} 
\Delta V(t)=\lambda \Delta I\left(1-e^{-t/\lambda } \right) 
\end{equation}

Figure \ref{fig:equil} shows a plot of \eqref{ZEqnNum190537} for 4 values of $\lambda$ (in minutes) with $\Delta I$=20 liters/minute. Note, for example, that for the $\lambda$ = 5 minute line the value at 30 minutes is 5 minutes x 20 liters/minute = 100 liters as predicted by \eqref{ZEqnNum190537}.


\begin{figure}[htbp]
    \centering
\includegraphics*[width=0.8 \textwidth]{approach_equil.png} 
\caption{$\Delta V$ vs. time from \eqref{ZEqnNum190537} for 4 different time constants $\lambda$=1/k (minutes)}
  \label{fig:equil}
\end{figure}




\section{Climate feedback}


 Feedback occurs when a disturbance in a system causes changes in the system which either enhance or suppress the original change.  If the changes that occur when the system is disturbed tend to increase the initial disturbance and destabilise the system, this is called an \textbf{amplifying feedback}.  Conversely, if the changes that occur when the system is disturbed tend to reduce the initial disturbance and stabilise the system, this is called a \textbf{stabilizing feedback}.

 One example of a stabilizing feedback is the bathtub outflow problem we solved in Section 1. Another example is a thermostat: a device that turns a heater on or off depending on the temperature.  When the temperature rises above a threshold level, the thermostat turns the heater off, reducing the amount of heating and lowering the temperature.  Conversely, when the temperature falls below the threshold level, the thermostat turns the heater on.  This ensures the room never gets too hot or too cold, keeping the room temperature stable

 An example of \textbf{amplifying feedback} occurs in a microphone connected to amplified speakers.  If the amplifier volume is too high, or if the microphone is brought too close to the speakers, small sounds picked up by the microphone will be made louder by the speakers. The amplified noise made by the speakers is then picked up by the microphone and amplified again and again, making the noise louder and louder until the speakers are making the loudest noise the amplifier is capable of emitting. 

\section{Stabilizing the climate with the black-body/Planck feedback}

\begin{figure}[htbp]
    \centering
\includegraphics*[width=0.8 \textwidth]{planck_curve.png}
\caption{Upward blackbody flux as a function of temperature}
  \label{fig:planck}
\end{figure}

The long-wave radiation emitted by an object is given by $I_{\uparrow } =-{\rm }\sigma T^{4} $, so as the temperature of the climate increases, the amount of long-wave radiation emitted by the Earth's surface (negative upward)  will increase and the impact on the TOA (top of atmosphere) radiative balance will be increasingly negative (cooling).  This increased cooling acts to reduce the temperature of the climate, causing a stabilising feedback.  This feedback is called the long-wave radiation feedback, or sometimes the Planck feedback.

 How large is this feedback?  Figure \ref{fig:planck} shows the rate of decrease of I${}_{G}$ with temperature.  From the graph we can estimate the slope at about 280 K as --(-450 -- (-250))/(300 -- 260) =-200 Wm${}^{-2}$/40 K = -5 Wm${}^{-2}$/K.  To find the exact value of the slope at 280 K use calculus:
 \begin{equation}
   \label{7)} 
\frac{dI_{G} }{dT} =\frac{d(-\sigma T^{4} )}{dT} =-4\sigma T^{3} =-4\sigma 280^{3} =-4.97\; Wm^{-2} K^{-1}  
\end{equation} 


 We can use the Planck feedback to get a better estimate of what happens as we heat the ocean with a constant forcing as we did in the Day 10 reading.  From Day 20 p. 5 we had the following equation:

 
 \begin{equation}
   \label{8)} 
\frac{d\Delta E}{dt} =\rho_{w} f_{w} D\frac{\Delta T}{dt} =\Delta F 
\end{equation} 
Now add in the cooling effect of the Planck feedback:
\begin{equation}
  \label{ZEqnNum244370} 
\frac{d\Delta E}{dt} =\rho _{w} f_{w} D\frac{d\Delta T}{dt} =\Delta F+\Delta R=\Delta F+\frac{dI_{G} }{dT} \Delta T=\Delta F-4\sigma T^{3} \Delta T 
\end{equation} 
Where we have defined $\Delta R$ as the \textbf{radiative response, }which in this case is the increase in the cooling flux due to the Planck feedback.  To solve \eqref{ZEqnNum244370}, define the following terms:

 The \textit{climate sensitivity}
 \begin{equation}
   \label{ZEqnNum400362} 
\lambda =\frac{1}{4\sigma T^{3} } {\rm \; K/(Wm}^{-2} ) 
\end{equation} 
The \textit{climate timescale}:
\begin{equation}
  \label{11)} 
\tau =\lambda \rho _{w} f_{w} D{\rm \; (seconds)} 
\end{equation} 
With these two definitions, we can rewrite \eqref{ZEqnNum244370} as:
\begin{equation}
  \label{ZEqnNum241455} 
\frac{\tau }{\lambda } \frac{\Delta T}{dt} =\Delta F+\Delta R=\Delta F-\frac{\Delta T}{\lambda }  
\end{equation} 
Comparing \eqref{ZEqnNum241455} with \eqref{ZEqnNum448161} we can see that they are nearly identical.  Convince yourself that the solution to \eqref{ZEqnNum241455} is given by:
\begin{equation}
  \label{ZEqnNum653852} 
\Delta T(t)=\lambda \Delta F\left(1-e^{-t/\tau } \right) 
\end{equation} 
When we go to the limit of long times $e^{-t/\tau}$ goes to 0 and $\Delta$T stops changing.  This gives us the definition of $\lambda$ as the \textbf{\textit{equilibrium climate sensitivity:}}
\begin{equation}
  \label{14)} 
\Delta T=\lambda \Delta F 
\end{equation} 
That is, larger values of $\lambda$ means larger final temperature increases for a given forcing -- i.e. high $\lambda$ climate systems are more sensitive to forcing.

\hfill \break


\textbf{Example Problem: }

$~$

Suppose an ocean layer with D=150 m is in equilibrium at a temperature of 280 K.  It undergoes a forcing increase of $\Delta F = 1\ W\,m^{-2}$.

$~$

\textbf{Part 1: If the only feedback is the Planck feedback, what is the final equilibrium temperature?}

$~$

First find the climate sensitivity due to the Planck feedback alone using equation
\eqref{ZEqnNum400362}.  At a temperature of 280 K, the Planck feedback is
$f_{pl}= -4 \sigma 280^{-4} = -5 \ W\,m^{-2}\,K^{-1}$ 

That means that the climate sensitivity due to the Planck feedback alone is

$\lambda = -1/f_{pl} = -1/(-5) = 0.2\ K/(Wm^{-2})$


 So ${\Delta T}_{final}=\lambda \Delta F = 0.2 \times 1$ =0.2 K

 $~$
 
\textbf{Part 2: What is the time constant $\tau$ for this problem?}

$~$

 $\tau=\lambda \rho_w f_w D$ = 0.2*1000*4218*150=4 years

\hfill \break
 

 The Planck feedback is often thought of as separate from the rest of the climate feedbacks, because it is a process which will occur in any climate---even the moon experiences this feedback---while things like cloud and ice feedbacks change when the amount of clouds and ice cover changes.  The other feedbacks are often thought of as modifying this basic feedback process.  This is the main difference between this reading and Dessler section 6.4 -- he considers the Planck feedback to be a basic part of the climate system, and doesn't count it in his equation 6.7 -- most modern work on climate feedback (including the problems in Assignment 4 uses the approach in this reading.

 \textbf{}


\section{Combining climate feedbacks}


 \textbf{To review -- Climate feedbacks} are changes to the net radiative flux into and out of the \textit{top of the Earth's atmosphere} in response to changes in the Earth's temperature. They can be included in the earth's energy budget equation in this form:
 \begin{equation}
   \label{ZEqnNum380732} 
\frac{d\Delta E}{dt} =\Delta F+\Delta R=\Delta F+f\Delta T 
\end{equation} 
Here, $\Delta E$ is the change in the column energy from a reference equilibrium value, the climate forcing $\Delta$\textit{F }(in W m${}^{-2}$) is a change in the net heat flux into the top of the Earth's atmosphere (positive downward) that is \textit{independent of temperature}, $\Delta$R is the planet's response to the forcing, which can also be written in terms of the climate feedback \textit{f} $\Delta$\textit{T} as a change in the TOA radiative flux that occurs as a result of a change in the Earth's temperature .  The \textbf{climate feedback factor} f\textit{ }has units of W m${}^{-2 }$K${}^{-1}$ and is the negative inverse of the climate sensitivity (i.e. \textit{f=}-1/$\lambda$).  If \textit{f} is positive, then the climate feedback is amplifying, and if f is negative then the climate feedback is stabilizing. For the climate to be stable, f\textit{ }must be negative; otherwise, changes in temperature would feed back on the forcing and drive the Earth to a completely new equilibrium.  However, the total climate feedback factor f  characterizes several different feedback processes, some of which are amplifying, and some of which are stabilizing.

\section{Major Climate feedbacks}

 In addition to the Planck feedback, there are many feedbacks that operate in the Earth's climate system, on many different time scales.  However, over the time scale of a human lifetime ($\sim$80 years) there are five main feedback processes that affect the climate.

 

\subsection{Water Vapor}

 Water vapour feedback arises because warm air can hold more water vapour than cooler air.  As the temperature of the climate system increases, the maximum amount of water the atmosphere can hold also increases.  Since most of the Earth surface is covered in ocean, this warming results in more evaporation and increases the amount of water vapour in the Earth's atmosphere.  Since water vapour is a powerful greenhouse gas, this increase in water vapour reduces the amount of long-wave flux  radiated from the atmosphere to space, causing the temperature to increase even more.  This makes water vapour feedback an amplifying feedback process.  It is positive because it acts to decrease the negative (upward) outgoing longwave radiation (OLW).

\subsection{Lapse rate}


 The strength of the greenhouse effect does not only depend on the amount of greenhouse gas in the atmosphere, but also on the atmospheric lapse rate=dT/dz.  If the lapse rate were zero, so that the top of the atmosphere were the same temperature as the bottom of the atmosphere, the top of the atmosphere would radiate energy to space at the same rate it absorbs energy from below and there would be no greenhouse effect.  As the lapse rate gets more negative, the top of the atmosphere becomes colder relative to the surface, and the greenhouse effect becomes stronger.

 Exactly how the lapse rate changes as the Earth's temperature changes is not obvious because it depends on many processes, such as the strength of convection and how quickly atmospheric water vapour declines with height.  However, simulations done using climate models indicate that as the Earth's temperature rises, the top of the troposphere tends to warm faster than the surface, making the lapse rate more positive.  If the lapse rate more positive, then the top-bottom temperature contrast and the greenhouse effect are reduced and more outgoing long-wave radiation escapes from the top of the atmosphere.  This makes the lapse rate feedback a stabilising feedback process, is negative because it increases OLW.

\subsection{Ice-albedo}


 Ice and snow generally reflect short-wave radiation much better than rock, dirt, or water.  Snow has an albedo of about $\alpha$ = 0.8-0.9 and ice has a slightly lower albedo between $\alpha$ = 0.5-0.7, while the ocean's albedo is around $\alpha$ = 0.06 and land is between $\alpha$ = 0.1-0.3.  As the Earth's temperature increases the ice and snow on mountains and near the poles begins to melt, reducing the Earth's albedo and causing more short-wave radiation to be absorbed by the surface, increasing the temperature further.  This process is called the ice-albedo feedback, and it is an amplifying feedback process. It has a positive sign because it increases the net positive downward shortwave.

\subsection{Clouds}

 The cloud feedback effect depends on several processes.  As we discussed in the Day 9 notes, low clouds tend to reflect more short-wave radiation, cooling the Earth, while high clouds allow short wave through but trap outgoing long-wave radiation, warming the Earth.  As the temperature of the Earth changes, the amount of high and low clouds will change, altering how much radiation clouds trap and reflect.  In addition, changes in the water vapour content of the atmosphere will alter the average size of the water droplets in the clouds, which will affect their albedo and emissivity.  Finally, changes in the atmospheric temperature will affect whether the clouds are made of water droplets or ice crystals, which will alter their albedo and emissivity as well.  

 The many ways in which clouds could change as the temperature changes makes the cloud feedback effect the most uncertain and least understood of all the climate feedbacks.  However, it appears that as the temperature of the Earth increases, the amount of low clouds decreases and the amount of high clouds increases, warming the climate further and making the cloud feedback an amplifying feedback process.


\section{Climate Sensitivity and Feedbacks}


 A primary goal of modern climate science is to understand and more accurately estimate the climate sensitivity, $\lambda$, which will tell us how much the CO${}_{2}$ emissions we are adding to the atmosphere will warm the climate.  Since the climate sensitivity is determined by the climate feedback factor, \textit{f }= -1/\textit{$\lambda$}, to calculate the sensitivity we must determine how climate feedbacks affect the energy budget of the Earth. 

 Repeating equation \eqref{ZEqnNum380732}:
 \begin{equation}
   \label{ZEqnNum704639} 
\frac{d\Delta E}{dt} =\Delta F+f\Delta T 
\end{equation} 
Where again, \textit{E} is the energy stored in the climate system in J m${}^{-2}$, $\Delta$\textit{F }is the climate forcing in W m${}^{-2}$, $\Delta$\textit{T }is the average temperature change of the climate (usually, but not always, taken to be the surface temperature), and \textit{f }is the climate feedback factor in W m${}^{-2}$ K${}^{-1}$.  This equation treats the various feedback processes as a single process, represented by \textit{f.}  In reality each feedback processes is studied separately by many different scientists, and to determine the total effect of all the climate feedbacks, we need to add all of the individual effects together.

 Separating the climate feedback factor into individual feedbacks is fairly easy to do.  We can approximate the effect of each of these feedbacks as a temperature-independent constant times the change in the Earth's temperature
 \begin{equation}
   \label{ZEqnNum851416} 
\frac{d\Delta E}{dt} =\Delta F+f_{PL} \Delta T+f_{WV} \Delta T+f_{A} \Delta T+f_{LR} \Delta T+f_{C} \Delta T 
\end{equation} 
Here, f${}_{pl }$is the feedback factor due to the Planck long-wave radiation feedback, f${}_{wv}$\textit{ }is that due to the water vapour feedback, \textit{f${}_{A}$}${}_{ }$is that due to the ice-albedo feedback, \textit{f${}_{LR}$}${}_{ }$is that due to the lapse rate feedback, and \textit{f${}_{C}$}${}_{ }$is that due to the cloud feedback (all in W m${}^{-2}$ K${}^{-1}$).  For simplicity, we can factor the temperature change $\Delta$\textit{T} out of each feedback term and write 
\begin{equation}
  \label{18)} 
\frac{d\Delta E}{dt} =\Delta F+\left(\sum f_{i}  \right)\Delta T 
\end{equation} 
Here $\Sigma$f${}_{i}$\textit{${}_{ }$}has replaced the \textit{f }in equation \eqref{ZEqnNum704639}, so the net climate feedback factor (f) is just the sum of each individual climate feedback factor.

 Since at equilibrium \textit{d$\mathit{\Delta}$E}/\textit{dt} = 0, we can write
 \begin{equation}
   \label{ZEqnNum249012} 
\Delta F=-\left(\sum f_{i}  \right)\Delta T 
\end{equation} 
but we already know from \eqref{ZEqnNum653852} that in equilibrium  $\Delta$\textit{T} = $\lambda$$\Delta$\textit{F}  and so the climate sensitivity can be calculated according to
\begin{equation}
  \label{ZEqnNum332906} 
\lambda =-\frac{1}{f} =-\frac{1}{\sum f_{i} }  
\end{equation} 
What is equation \eqref{ZEqnNum332906} saying, in words?  As long as the individual feedbacks sum up to a negative number, then equation \eqref{ZEqnNum249012} says that there can be a positive temperature change $\Delta$T  that would counteract a positive forcing $\Delta$F (and this also works in the opposite direction, such that a negative cooling $\Delta$F can be offset by a negative temperature change.  If the feedbacks sum to a positive number, then there is nothing that can hold back dE/dt from growing indefinitely for positive $\Delta$F.   This situation where the net amplifying feedback is amplifying is called ``runaway feedback''.$ $ 


\section{Climate Feedback Factors}


 Now the question becomes how to calculate the climate feedback factors.  The climate feedback factor is defined as the change in the top-of-atmosphere (TOA) heat flux due to a change in the Earth's temperature.  Therefore, we can calculate the climate feedback factor by calculating how much the TOA heat flux changes when the temperature changes.  It also makes it simpler to show how the various feedback factors are calculated below.

 As the surface temperature changes, various components of the climate (cloud cover, lapse rate, water vapour, etc) will also change, producing a change in the radiation at the top of the atmosphere we will denote as $\Delta$R. Whatever the component, we can then (using the chain rule) calculate the feedback factor 
 \begin{equation}
   \label{21)} 
f_{x} =\frac{\Delta R}{\Delta T} =\left(\frac{\Delta R}{\Delta climate_{x} } \right)\left(\frac{\Delta climate_{x} }{\Delta T} \right) 
\end{equation} 
Where the subscript x is either clouds, water vapour, sea ice, etc. Some processes are easier to characterize than others.  The long-wave radiation emitted by the Earth depends on temperature directly, via $I_{\uparrow } =-{\rm }\sigma T^{4} $, so we calculate the Planck feedback factor by calculating
\begin{equation}
  \label{ZEqnNum227245} 
f_{PL} =\frac{\Delta R}{\Delta T} =\frac{dI_{\uparrow } }{dT} =-{\rm }\sigma 4T^{3}  
\end{equation} 
(note the assumption that $\epsilon$ isn't changing -- we're tracking the emissivity changes in other feedbacks like H${}_{2}$O. The true calculation is complicated by having to average these values over the whole surface of the Earth.  However, doing so results in an estimate of $f_{PL}$= - 3.2 W $m^{-2}\, K^{-1}$, which is very similar to the feedback one calculates for $\varepsilon$ = 1 and T = 240 K (the temperature near the top of the troposphere) using Equation \eqref{ZEqnNum227245}, i.e$-4\sigma 240^{3} =-3.1\; W\, m^{-2} \, K^{-1} $ 

 The water vapour feedback is given by
 \begin{equation}
   \label{23)} 
f_{WV} =\frac{\Delta R_{WV} }{\Delta T} =\left(\frac{\Delta R}{\Delta H_{2} O} \right)\left(\frac{\Delta H_{2} O}{\Delta T} \right) 
\end{equation} 
where the \textit{H}${}_{2}$\textit{O }represents the water vapour content of the Earth's atmosphere.  The ice-albedo feedback is given by
\begin{equation}
  \label{24)} 
f_{A} =\frac{\Delta R_{A} }{\Delta T} =\left(\frac{\Delta R}{\Delta {\rm ice\; area}} \right)\left(\frac{\Delta {\rm ice\; area}}{\Delta T} \right) 
\end{equation} 
where ``ice area'' is the amount of the Earth's surface covered in ice.  The lapse rate feedback is given by
\begin{equation}
  \label{25)} 
f_{WV} =\frac{\Delta R_{LR} }{\Delta T} =\left(\frac{\Delta R}{\Delta LR} \right)\left(\frac{\Delta LR}{\Delta T} \right) 
\end{equation} 
Finally, the cloud feedback is given by
\begin{equation}
  \label{26)} 
f_{WV} =\frac{\Delta R_{C} }{\Delta T} =\left(\frac{\Delta R}{\Delta {\rm clouds}} \right)\left(\frac{\Delta {\rm clouds}}{\Delta T} \right) 
\end{equation} 
where ``clouds'' represents changes in the amount, average height, and properties of the clouds in a cloud model and $\Delta$R${}_{c}$ represents the change in radiation at the top of the atmosphere due to the change in clouds calculated using a radiation code like Modtran.

\begin{figure}[htbp]
    \centering
\includegraphics*[width=0.8 \textwidth]{feedback_factors.png} 
\caption{Feedback Magnitudes, not including the Planck feedback. (IPCC AR4 2007)}
  \label{fig:factors}
\end{figure}


 Figure \ref{fig:factors} shows estimates for WV=water vapour, C = cloud, A=Albedo, LR=lapse rate, WV+LR=combined water vapour and lapse rate, and ALL=everything calculated using climate models for each of the climate feedback factors, not including the Planck feedback.  Because each model makes different approximations in order to simulate atmospheric physics, each model has slightly different estimates of the climate feedback strengths.  The water vapour feedback has the largest magnitude, with a feedback factor of (1.80 $\pm$ 0.18) W m${}^{-2}$ K${}^{-1}$.  The cloud feedback C is the most uncertain feedback, with a feedback factor of  (0.69 $\pm$ 0.38) W m${}^{-2}$ K${}^{-1}$.  Finally, the albedo feedback has a feedback factor of (0.26 $\pm$ 0.08) W m${}^{-2}$ K${}^{-1}$, and the lapse rate feedback has a feedback factor of (-0.84 $\pm$ 0.26) W m${}^{-2}$ K${}^{-1}$.  When the water vapour and lapse rate feedbacks are added together (WV + LR) there is cancellation, because models with active convection have both strong negative lapse rate feedback (because clouds reduce the temperature contrast between the upper atmosphere and the surface) and strong positive water vapour feedback (because clouds move water vapour up to higher altitudes where it has a larger greenhouse effect).   The best estimate of the net climate feedback factor from these four climate feedbacks (ALL in figure 1) is approximately (1.91 $\pm$ 0.2) W m${}^{-2}$ K${}^{-1}$.  Combining this with the Planck feedback value of -3.2 W m${}^{-2}$ K${}^{-1}$ gives a total climate feedback of \textit{c} = -1.29 W m${}^{-2}$ K${}^{-1}$, and a climate sensitivity of  $\lambda$ = 0.78 K (W m${}^{-2}$)${}^{-1}$.

 Note that the standard deviation in the total feedback value (0.2 W m${}^{-2}$ K${}^{-1}$) is \textit{lower} than the cloud feedback uncertainty (0.38 W m${}^{-2}$ K${}^{-1}$) or the lapse rate feedback uncertainty (0.26 W m${}^{-2}$ K${}^{-1}$).  This is because we can measure the \textit{net} climate feedback using paleoclimate data much better than we understand each individual feedback process.  Each model is designed so its net climate feedback matches our observations of climate feedback, but each does so in a different way, so that models with a strong cloud feedback may have a weaker water vapour feedback or a stronger lapse rate feedback.  This is a serious problem which limits the usefulness of \textbf{\textit{regional}} climate model predictions, since the climate change in any given region experiences will depend on the details of how the feedbacks adjust the cloud cover, water vapour, lapse rate, and surface albedo.  In other words, we are confident the Earth is warming, but we are unsure exactly how this warming will change the climate.


\section{Feedback Changes}


 In writing equation \eqref{ZEqnNum851416} we made the implicit assumption that the strength of the climate feedbacks don't change as the temperature changes.  In reality, the strength of each climate feedback will change as the climate alters.  With large temperature changes, some of the feedback processes can even disappear completely; for example, if all the ice cover on the Earth were to melt, the ice-albedo feedback would stop.  Conversely, it is possible for temperature changes to start new feedback processes; an example of this is the concern that methane locked in the Arctic permafrost will be released as the Arctic warms, creating a new positive methane feedback.  As the climate feedbacks change, the climate sensitivity also changes; these changes can be calculated using equation \eqref{ZEqnNum332906}.

 If the sum of the climate feedback factors were ever to achieve a net positive value, this would trigger a situation known as \textbf{runaway feedback}: an increase in the temperature of the climate would increase the net heat flux into the climate system and the temperature of the climate would increase indefinitely.  This process would stop only when the climate had changed to such a degree that the net climate feedback became negative once more.  

 Runaway feedback events have previously occurred in Earth's history.  Several times in Earth's past, before multicellular life appeared, runaway ice-albedo feedback resulted in the entire Earth being covered in ice, a climate state referred to as ``Snowball Earth''.  Once the entire Earth had frozen, the ice-albedo feedback stopped, as the Earth's albedo could not increase any further.  On the other extreme, 55 million years ago during the Eocene, shortly after the extinction of the dinosaurs, a massive increase in the amount of CO${}_{2}$ in the atmosphere resulted in global temperatures increasing by 5 K over about 10,000 years.  It has been hypothesized that this temperature increase resulted from a positive feedback caused by the release of methane clathrate deposits on the ocean floor.  As the temperature of the Earth increased, methane clathrates released into the atmosphere trapped more heat, further increasing the temperature of the Earth and causing even more methane to be released.  This feedback would have stopped once all the methane deposits were released to the atmosphere.  There are also examples of runaway feedback on other planets; because Venus is closer to the sun than the Earth is, Venus experienced a runaway water vapour feedback early in its history resulting in surface temperatures well over 700 K.  This feedback stopped only once all of Venus' oceans had evaporated into water vapour.

 

\section{Summary}

\begin{itemize}
\item \textbf{\textit{ }}Climate Feedbacks are processes that depend on the temperature of the Earth and alter the radiative balance of the Earth.

\item  an \textbf{amplifying feedback} system responds to a disturbance by increasing the initial disturbance, destabilising the system.  A \textbf{stabilising feedback} system responds to a disturbance by decreasing the initial disturbance, stabilising the system.  amplifying and stabilising feedbacks are not good or bad; ``amplifying'' and ``negative'' refer to whether the feedback has the same or the opposite sign as the forcing that disturbs the system.

\item  The five major short term climate feedbacks come from blackbody radiation, water vapour, the atmospheric lapse rate, the ice-albedo feedback, and clouds.

\item  Black body or Planck feedback occurs when a temperature change alters the outgoing long-wave radiation emitted from the top of the atmosphere.  This is a stabilising feedback.

\item  Water vapour feedback occurs when a temperature change alters the atmospheric water vapour content, altering the strength of the greenhouse effect.  This is an amplifying feedback.

\item  Lapse rate feedback occurs when a temperature change alters the atmospheric lapse rate, altering the strength of the greenhouse effect.  This is a stabilising feedback.

\item  Ice-albedo feedback occurs when a temperature change alters the amount of ice on Earth, changing the Earth's albedo.  This is an amplifying feedback.

\item  Cloud feedback occurs when a temperature change alters the average amount of cloud cover, changing the amount of long wave trapped and short wave reflected by the clouds. This is an amplifying feedback. 

\item  The strength of climate feedbacks determine the size of the climate sensitivity. 

\item  $$\lambda =-\frac{1}{\sum f_{i} } $$

\item  The strength of the climate feedbacks can be determined by calculating how much the climate changes when the feedback system changes and multiplying that by how much the heat flux at the top of the atmosphere changes when the feedback system changes. 

\item  $f_{x} =\frac{\Delta R}{\Delta T} =\left(\frac{\Delta R}{\Delta climate_{x} } \right)\left(\frac{\Delta climate_{x} }{\Delta T} \right)$

\item  The water vapour feedback is the strongest climate feedback, and the cloud feedback is the most uncertain.

\item  If the net climate feedback factor ever goes above 0, the climate would enter into \textbf{runaway feedback}, in which temperature increases would increase the amount of heat flux into the climate, and the climate would change radically.

\end{itemize}


 


\section{Further Reading:}

Chapter 6 of Dessler, A. E., ``Introduction to modern climate change'' (2016), Cambridge University Press.




\end{document}


%%% Local Variables:
%%% mode: latex
%%% TeX-master: t
%%% End:
